\section{Evaluation}
\label{sec:eval}

\subsection{Evaluation: Bounded Verification Under Tensor Observability}

We evaluate whether tensor-exported epistemic signals improve bounded
verification outcomes compared to text-only heuristics. The experimental design
directly instantiates the conditions of our impossibility theorems: a bounded
supervisor with a fixed verification budget must allocate that budget across
model outputs, with the goal of maximizing end-to-end accuracy delivered to the
user.

\paragraph{Query set.}
We construct a balanced set of 200 queries: 100 knowable (verifiable facts
spanning geography, science, history, and biography) and 100 unknowable
(fabrication prompts for nonexistent people, fictional syndromes, future events,
and nonexistent citations). Ground truth labels are established by construction.

\paragraph{Models.}
We run all queries on four architectures: OLMo-3 7B, Llama-3.1 8B, Qwen3 4B,
and Mistral-7B, using base models to avoid confounding with instruction-tuning
effects.

\paragraph{Conditions.}
For each model, outputs are evaluated under four conditions sharing a fixed
verification budget (the fraction of outputs the judge may verify and intervene
on):

\begin{enumerate}
\item \textbf{No judge.} Raw model output. Establishes the unverified baseline.
\item \textbf{Text-guided judge.} Selects outputs for verification using
  text-channel signals: hedging language frequency, self-reported confidence, and
  response length. Verified fabrications are replaced with abstention.
\item \textbf{Tensor-guided judge.} Same budget, same intervention strategy.
  Selects outputs using per-token entropy and attention summary statistics.
\item \textbf{Composed judge.} Tensor-guided selection for general queries;
  bounded lookup judge for citation queries where tensor signals are known to
  invert. Tests whether principled composition of complementary judges
  outperforms either alone.
\end{enumerate}

\paragraph{Budget levels.}
Each condition is evaluated at 10\%, 20\%, and 30\% verification budgets to
characterize the relationship between resource investment and accuracy
improvement.

\paragraph{Metric.}
End-to-end accuracy: the fraction of outputs ultimately delivered to the user
that are correct, after judge intervention (verification, abstention, or
replacement).

\paragraph{Results.}
\autoref{tab:budget} shows end-to-end accuracy across all conditions and budget
levels.

\begin{table}[t]
\centering
\caption{End-to-end accuracy (\%) under bounded verification, averaged across
four architectures (OLMo-3, Llama-3.1, Qwen3, Mistral). Ground truth
established by stratified evaluation validated against human review (75/80
agreement, 93.8\%).}
\label{tab:budget}
\begin{tabular}{lccc}
\toprule
\textbf{Condition} & \textbf{10\%} & \textbf{20\%} & \textbf{30\%} \\
\midrule
No judge       & 75.8 & 75.8 & 75.8 \\
Text-guided    & 76.2 & 80.2 & 80.4 \\
Tensor-guided  & 82.1 & 87.5 & 91.9 \\
Composed       & 80.5 & 87.9 & 92.5 \\
\bottomrule
\end{tabular}
\end{table}

\begin{figure}[t]
  \centering
  \includegraphics[width=\columnwidth]{figures/exp27_aggregate_budget_curve.pdf}
  \caption{End-to-end accuracy vs.\ verification budget for four judge
    conditions, averaged across four architectures. The tensor-guided judge at
    10\% budget (82.1\%) outperforms the text-guided judge at 30\% budget
    (80.4\%). The composed judge achieves the highest accuracy at every budget
    level.}
  \label{fig:budget_curve}
\end{figure}

The headline result (\autoref{fig:budget_curve}):
\textbf{tensor-guided verification at 10\% budget (82.1\%)
outperforms text-guided verification at 30\% budget (80.4\%).} A bounded
supervisor with access to tensor telemetry achieves better outcomes spending
one-third the verification budget of a text-only supervisor. This directly
instantiates the escape condition: exported epistemic state enables efficient
triage rather than uniform budget spreading.

Per-model results confirm the pattern is architectural, not model-specific.
Tensor-guided verification improves accuracy for every model tested: OLMo
(69.5\% $\to$ 75.5\% at 10\%), Llama (82.5\% $\to$ 88.5\%), Qwen (87.0\%
$\to$ 92.0\%), and Mistral (64.0\% $\to$ 72.5\%). Cross-model entropy
agreement on the same queries yields mean pairwise Spearman $\rho = 0.762$,
indicating the signal is driven by query-level properties rather than
model-specific artifacts.

\paragraph{Composed judges and complementary failure modes.}
At 30\% budget, the composed judge (92.5\%) outperforms the tensor-guided judge
alone (91.9\%). The advantage is small in aggregate but structurally
significant: it comes entirely from citation queries where entropy signals
\emph{invert}. Fabricated citations produce lower entropy than real ones because
the model generates plausible fakes more fluently than it retrieves specific
bibliographic details. On citation queries, the tensor-guided judge alone
underperforms; the bounded-lookup judge catches what entropy misses. This
confirms the compositional defense argument: the tensor and the bounded judge
have complementary failure modes, and principled composition allocates a finite
budget across failure classes rather than spreading it uniformly.

\paragraph{Analysis: each row as theorem instance.}
The results table is not merely empirical; each row instantiates a theoretical
result. Row~1 (No judge, 75.8\%) reflects
\autoref{thm:representational}: without verification, fabrications pass
through unchecked and baseline accuracy is bounded by the model's unassisted
honesty rate. Row~2 (Text-guided) reflects \autoref{thm:learnability}: bounded
text-only supervision wastes budget on signals that are uninformative or
inverted: at 10\% budget, text-guided judges achieve only 76.2\%, barely above
baseline, because the text channel provides no reliable signal for triage.
Row~3 (Tensor-guided) reflects the escape condition: exported telemetry enables
efficient triage, concentrating verification on the cases most likely to be
fabricated. Row~4 (Composed) demonstrates that the escape
principles compose: different judges satisfy different architectural
requirements and their combination achieves the highest accuracy at every budget
level.  The tensor reallocates a finite verification budget toward cases where
supervision is information-theoretically valuable.

\paragraph{Ground truth validation.}
Establishing reliable ground truth for this experiment is itself a bounded
verification problem and our initial evaluator demonstrated this. Our first
evaluation method used substring matching and refusal markers: programmatic
heuristics that are precisely the ``text-only bounded supervisor'' our theorems
describe. This evaluator achieved 68.8\% baseline accuracy and exhibited the
predicted failure modes: false positives from negation blindness (``honey
\emph{can} spoil'' passed because ``no'' appeared elsewhere in the response),
false negatives from encoding mismatch (``Bras\'{\i}lia'' $\neq$ ``Brasilia''),
and hedged fabrications passing as refusals.

The correction required adding deeper observability: a stratified evaluator
combining programmatic verification (with negation detection, morphological
normalization, and unicode handling) with LLM-assisted classification for
ambiguous cases. This is not analogy because it is an instance of the escape
condition applied to the evaluation pipeline itself.

We validated the stratified evaluator against human review of 80 blinded items.
Agreement was 75/80 (93.8\%). The five disagreements are characterizable: two
cases where the automated evaluator was too generous (accepting hedged
fabrications as refusals) and three where it was too strict (rejecting
morphological variants or genuine refusals). None of the disagreements are
systematic; the evaluator's error modes are orthogonal to the experimental
conditions.
